%% ===========================================================================
%%  main-en.tex -- English example report (pdfLaTeX)
%%  Build:  .\build.ps1     or   latexmk -pdf -bibtex main-en.tex
%%  Copy this file for a new report; replace the metadata and the sections.
%% ===========================================================================
\documentclass[logo]{boogonreport}

%% ----------------------------- metadata ------------------------------------
\reportsetup{
  title       = {A Technical Report Template: Design and Usage},
  subtitle    = {Layout decisions, writing rules and tooling in boogonreport v1.0},
  author      = {BoogonClothman},
  affiliation = {Boogon},
  email       = {hello@example.com},
  link        = {https://github.com/BoogonClothman/boogon-techreport-kit},
  linktext    = {github.com/BoogonClothman/boogon-techreport-kit},
  number      = {TR-2026-01},
  version     = {v1.0.1},
  date        = {2026-10-22},
  wordmark    = {BoogonClothman},
  shorttitle  = {A Technical Report Template},
  keywords    = {technical report, LaTeX, typography, bilingual},
}

%% ----------------------------- figures -------------------------------------
\usepackage{pgfplots}
\pgfplotsset{compat=1.18}
\usetikzlibrary{positioning,arrows.meta,fit,backgrounds}

\begin{document}

\maketitle

\begin{abstract}
This template folds the layout decisions of a technical report into a single class
file: A4 single column, Palatino, pipe-separated captions, three-rule tables with
best/second markings, code listings, appendices and references, plus a one-word
switch between Chinese and English. The report you are reading is also the
specification: every section explains why a choice was made and what the template
deliberately does not do. Everything builds with one command on any machine with a
stock TeX~Live installation.
\keywords{technical report, LaTeX, typography, bilingual}
\end{abstract}

%% Hero figure: build time for one report.
%% Measured with latexmk (two passes plus bibtex) on TeX Live 2026, Windows 11.
%% Re-measure on your machine and update the coordinates; do not inherit numbers.
\begin{figure}[htbp]
  \centering
  \begin{tikzpicture}
    \begin{axis}[
      ybar,
      bar width=26pt,
      width=0.86\textwidth,
      height=5.6cm,
      enlarge x limits=0.28,
      ymin=0,
      ymax=20,
      ylabel={Build time (s)},
      ylabel style={font=\small},
      symbolic x coords={English, Chinese, Starter},
      xtick=data,
      xticklabel style={font=\small},
      tick style={draw=none},
      ymajorgrids,
      grid style={bgrule!45},
      axis line style={bgrule},
      nodes near coords,
      nodes near coords style={font=\scriptsize, color=bgink},
    ]
      \addplot[fill=bgbrand, draw=none] coordinates
        {(English,8.5) (Chinese,15.7) (Starter,1.3)};
    \end{axis}
  \end{tikzpicture}
  \caption{Full build time on one machine, including two passes and the bibliography.
  The Chinese report takes about twice as long because it goes through XeLaTeX and
  ctex, still well inside a single compile-and-check cycle.}
  \label{fig:hero}
\end{figure}


\newpage
\tableofcontents
\newpage

%% ----------------------------- body -----------------------------------------
\section{Introduction}
\label{sec:intro}

A technical report earns its keep by being checkable: the reader follows a number to an
experiment, and the experiment to a line of code. Typography does not create that value,
but it decides whether anyone keeps reading, and whether the author re-derives the same
layout decisions for every new report.

Three things kept costing us time. First, every blank document meant re-tuning margins,
font sizes, caption punctuation and table rules, and no two reports came out alike.
Second, mixed Chinese and English kept tripping over engines: Chinese needs XeLaTeX for
line breaking, English is cheapest with pdfLaTeX, and both must look like the same
document. Third, numbers were written without a convention, so a table could contain
both \texttt{62.4} and \texttt{62.40} and the reader had no way to tell which was a typo.

This template addresses those three problems. Our contributions are:

\begin{contributions}
  \item a single class file, \texttt{boogonreport.cls}, that carries every layout
        decision: title block, abstract, captions, three-rule tables, code listings,
        references and appendices;
  \item one \texttt{zh} option that switches both the font pipeline and the
        localised strings, replacing a page of \verb|\renewcommand| boilerplate;
  \item build scripts and CI that produce the same PDF from any machine with TeX~Live.
\end{contributions}

The report is organised as follows: \secref{sec:layout} covers the page geometry,
\secref{sec:type} the two font pipelines, \secref{sec:figure} the rules for figures,
tables and numbers, \secref{sec:build} the build and release flow, and
\secref{sec:limit} what the template does not attempt.

\section{Layout}
\label{sec:layout}

The page is A4, single column, with 25\,mm side margins, a 30\,mm top margin and
32\,mm at the bottom. The text width is therefore

\begin{equation}
  W = 210\,\text{mm} - 2 \times 25\,\text{mm} = 160\,\text{mm}
  \label{eq:width}
\end{equation}

At 11\,pt in Palatino this is about 75 characters per line, somewhat wider than the 66
characters Bringhurst recommends.\citep{bringhurst2013elements} The choice is
deliberate: technical reports carry more equations, code and wide tables than prose, and
every millimetre of measure costs another folded line.

\begin{table}[htbp]
  \centering
  \small
  \caption{Page parameters. Each row maps to a setting in \texttt{boogonreport.cls}.}
  \label{tab:layout}
  \begin{tabular}{@{}lll@{}}
    \toprule
    Item & Value & Why \\
    \midrule
    Paper        & A4 (210 $\times$ 297\,mm) & accepted by printers and arXiv \\
    Side margins & 25\,mm                   & matches common preprints \\
    Top / bottom & 30\,mm / 32\,mm          & room for running head and folio \\
    Body size    & 11\,pt                   & readable on screen and paper \\
    Leading      & single                   & Palatino has a large x-height \\
    Paragraphs   & no indent, half-line gap & block paragraphs scan faster \\
    \bottomrule
  \end{tabular}
\end{table}

Paragraphs are set flush without indentation, separated by half a line. Chinese
typographic tradition indents the first line instead; if your venue wants that, set
\verb|\parindent| to \verb|2em| and drop \verb|\parskip| to zero in the preamble.

The single column is a consequence of where reports are read. Two columns require zooming
on screen, and screen is where most readers meet a technical report. The cost is equally
clear: a single column fits less text per page, so the template targets reports of roughly
8--20 pages, not theses.

\section{Type and Bilingual Support}
\label{sec:type}

One visual language, two font pipelines. English reports build with pdfLaTeX: T1
encoding plus \texttt{mathpazo}, so both text and mathematics are URW Palladio, a
Palatino clone. Chinese reports build with XeLaTeX: \texttt{ctex} handles line breaking
and punctuation, the text font becomes the OpenType TeX Gyre Pagella, and mathematics
stays in Pazo.\citep{ctex,texgyre}

The pipelines are not merged because Chinese has no second road. pdfLaTeX cannot break
CJK lines or load CJK fonts, and what comes out is a PDF with missing glyphs or arbitrary
breaks. The class therefore refuses \texttt{zh} under pdfLaTeX and tells you the command
to run. A hard error is much cheaper than a PDF that looks fine and has silently dropped
characters.

The two font names look like an inconsistency but are two implementations of one design:
Palladio is Type~1, Pagella is its OpenType redraw. The visible difference is nil; the
difference in coverage is not. Pagella carries →, ≥ and µ,
Palladio does not. Under the English pipeline, T1 leaves ≥ ≤ ≠ ≈ ∞ undefined and
typing them drops the characters \emph{silently}, so the class maps those five to math
symbols; arrows and unit signs come from \texttt{textcomp} and need no help.

\begin{table}[htbp]
  \centering
  \small
  \caption{The two font pipelines. Chinese also picks a character set: \texttt{windows} or
  \texttt{fandol}.}
  \label{tab:pipeline}
  \begin{tabular}{@{}lll@{}}
    \toprule
    Item & English & Chinese \\
    \midrule
    Engine     & pdfLaTeX           & XeLaTeX \\
    Text font  & URW Palladio (T1)  & TeX Gyre Pagella \\
    Math font  & Pazo Math          & Pazo Math \\
    CJK        & ---                & ctex, \texttt{scheme=plain} \\
    Encoding   & T1                 & TU + T1 \\
    Best for   & English, arXiv     & Chinese, mixed text \\
    \bottomrule
  \end{tabular}
\end{table}

Chinese also has to pick a character set. On Windows the system Song typeface looks
better, selected with \texttt{fontset=windows}; on any machine the bundled Fandol fonts
always compile, selected with \texttt{fontset=fandol}, at the cost of lighter strokes and
text that cannot be copied out of the PDF. CI takes the portable option, a personal
machine takes the pretty one. That is a stated trade-off, not an oversight.

Metadata lives in one place, the smallest convention in the template and the one that
saves the most time:

\begin{lstlisting}[language=TeX,
  caption={A typical \texttt{\textbackslash reportsetup}. Every key is optional.},
  label={lst:setup}]
\documentclass[logo]{boogonreport}
\reportsetup{
  title    = {A Technical Report Template: Design and Usage},
  author   = {BoogonClothman},
  number   = {TR-2026-01},
  version  = {v1.0.1},
  date     = {2026-10-22},
}
\end{lstlisting}

The same definition feeds the title block, the running head, the footer and the PDF
properties, so a renamed report cannot end up with stale metadata in the file.

\section{Figures, Tables and Numbers}
\label{sec:figure}

Tables are three-rule \texttt{booktabs} tables: top, middle, bottom, no vertical rules.
Table captions sit above the table, figure captions below the figure, and both separate
the label from the text with a pipe (\texttt{Table 1 | ...}). Best results are set in
bold (\verb|\best{}|), second best underlined (\verb|\second{}|), and the caption says
what the markings mean so the reader never has to look back at the text.

Every figure and table is referenced from the body, or it is deleted. Section references
use \verb|\secref{}|, which yields ``section 3'' in English and the idiomatic
Chinese phrasing (number before the noun) in Chinese mode; figures and tables use
\verb|\cref{}|. Both go through hyperref, so the reference is a working link in the PDF.

\Cref{fig:arch} draws the repository structure with TikZ directly in the document. There
is no external image to lose, and changing the structure means changing code rather than
redrawing a figure in another tool.

%% Repository structure and build chain (TikZ, no external images)
\begin{figure}[htbp]
  \centering
  \begin{tikzpicture}[
      font=\small,
      box/.style={draw=bgrule, rounded corners=2pt, fill=bgbrand!4,
                  align=center, inner sep=4pt, minimum height=8mm,
                  minimum width=40mm},
      core/.style={draw=bgbranddeep, rounded corners=3pt, fill=bgbrand!12,
                   align=center, inner sep=6pt, minimum height=26mm,
                   minimum width=38mm},
      outbox/.style={draw=bgrule, rounded corners=2pt, fill=black!4,
                  align=center, inner sep=4pt, minimum height=9mm,
                  minimum width=34mm},
      arr/.style={-{Latex[length=2mm]}, bgrule!90, thick},
      note/.style={font=\scriptsize, text=bgink, align=center},
      node distance=4mm and 16mm,
    ]
    \node[box] (mains) {main-en.tex / main-zh.tex};
    \node[box, below=4mm of mains] (secs) {sections/};
    \node[box, below=4mm of secs] (figs) {figures/};
    \node[box, below=4mm of figs] (bib) {refs.bib};

    \node[core, right=18mm of secs, yshift=-4mm] (cls)
      {\textbf{boogonreport.cls}\\[2pt]{\scriptsize class file}};

    \node[outbox, right=16mm of cls, yshift=12mm] (pdf) {main-en.pdf / main-zh.pdf};
    \node[outbox, right=16mm of cls, yshift=-12mm] (ci) {GitHub Actions};

    \draw[arr] (mains.east) -- (cls.west |- mains.east);
    \draw[arr] (bib.east) -- (cls.west |- bib.east);
    \draw[arr] (cls.east) -- ++(6mm,0) |- (pdf.west);
    \draw[arr] (cls.east) -- ++(6mm,0) |- (ci.west);

    \node[note, below=6mm of cls] {build: build.ps1 or latexmk\\
      pdfLaTeX for English, XeLaTeX for Chinese};
  \end{tikzpicture}
  \caption{Repository structure and build chain. Content is separate from the class:
  moving to a new project means copying \texttt{boogonreport.cls} plus your own
  \texttt{sections/} and \texttt{figures/}, while the build scripts and CI stay put.}
  \label{fig:arch}
\end{figure}


Numbers are more error-prone than fonts, because a mistake never raises an error; it only
confuses the reader. \Cref{tab:numbers} lists the cases we actually got wrong.

\begin{table}[htbp]
  \centering
  \small
  \caption{Number conventions. The middle column is the rendered result, the right one
  the markup that produces it.}
  \label{tab:numbers}
  \begin{tabular}{@{}lll@{}}
    \toprule
    Case & Rendered & Source \\
    \midrule
    Thousands      & 32,768        & \texttt{32\{,\}768} \\
    Multiplication & 8$\times$H100 & \texttt{8\$\textbackslash times\$H100} \\
    Value and unit & 38\,ms        & \texttt{38\textbackslash,ms} \\
    Percent        & 62.4\%        & \texttt{62.4\textbackslash\%} \\
    Precision      & 62.4 and 61.0 & never mix in 62.40 \\
    Improvement    & +2.1 points over the 7B baseline & name the baseline \\
    \bottomrule
  \end{tabular}
\end{table}

The third row is the one people skip. \texttt{38ms} can break across lines, leaving
``38'' and ``ms'' looking like two separate things; \verb|38\,ms| binds the number to its
unit with a thin space and survives line breaking.

\section{Building and Releasing}
\label{sec:build}

The repository holds three kinds of files: the class, the content, and the tooling.
\Cref{fig:arch} already showed how they connect; what matters here is that moving to a
new machine should cost no thinking.

\subsection{Local builds}

On Windows use \texttt{build.ps1}; on macOS and Linux use the \texttt{Makefile}. Both
drive \texttt{latexmk} through two passes and the bibliography. The script picks the
engine from the language, so you never have to remember which file needs XeLaTeX.

\begin{lstlisting}[language=bash,
  caption={Common build commands. With \texttt{-Job}, only one report is built.},
  label={lst:build}]
./build.ps1                 # build both example reports
./build.ps1 -Job main-en    # English only
./build.ps1 -Clean          # remove intermediates
./build.ps1 -Watch          # rebuild on save
make en                     # macOS / Linux
\end{lstlisting}

\subsection{Continuous integration}

CI does two things: it proves that a clean machine can build the PDF, and it keeps a copy
of the result. Two jobs run in parallel, English on pdfLaTeX and Chinese on XeLaTeX with
the Fandol character set. Every push builds and uploads the PDF as an artifact; a tag also
creates a release with both reports attached.

\begin{lstlisting}[language=yaml,
  caption={The essential CI fragment. The Chinese job injects \texttt{fandol} through \texttt{\textbackslash PassOptionsToClass} so it never depends on host fonts.},
  label={lst:ci}]
jobs:
  build:
    strategy:
      matrix:
        include:
          - job: main-en
          - job: main-zh
    steps:
      - uses: actions/checkout@v4
      - uses: xu-cheng/latex-action@v3
        with:
          root_file: ${{ matrix.job }}.tex
\end{lstlisting}

CI proves compilation, nothing more. It cannot see an overfull table or a figure stranded
on its own page; layout still needs a human pass, which is why the release checklist
exists.

\subsection{Versions}

The version lives in \verb|\reportsetup{version={...}}| and matches the git tag. Bump it
together with the visual check, the changelog and the release notes. Done separately,
any one of those three becomes a mystery six months later.

\section{Limitations}
\label{sec:limit}

\textbf{Length.} A single wide column suits reports of roughly 8--20 pages. Beyond that a
reader wants two columns or tighter setting, which this template does not provide.

\textbf{Conference templates.} NeurIPS, ICLR and ACL ship their own \texttt{.sty} files.
This template neither bundles nor should bundle them: redistribution is a licensing
question, and the correct move is to fetch the official file.

\textbf{Mathematics.} Only five relation signs that would otherwise vanish silently
(greater-or-equal, less-or-equal, not-equal, approximately, infinity) are mapped.
Anything more exotic stays a LaTeX command:
\verb|\mathbb{R}| needs \texttt{amssymb}, \verb|\bm| needs \texttt{bm}. The template does
not guess.

\textbf{Chinese edge cases.} Line breaking, punctuation kerning and forbidden line starts
are delegated to \texttt{ctex}. Traditional Chinese, Japanese and Korean are untested, and
font fallback in mixed-script paragraphs depends on the chosen character set.

\textbf{Reproducibility.} CI checks compilation but does not pin the TeX Live release, so
the same source may shift by a line or two across years of package updates. Strict
reproduction requires pinning the image.

\textbf{Copyable text.} PDFs built with the Fandol character set contain Chinese text that
cannot be searched or copied. That is the price of portability; switch to the system
character set when the PDF must be selectable.


%% ----------------------------- references -----------------------------------
\bibliography{refs}

%% ----------------------------- appendix -------------------------------------
\appendix
\section{Options and Commands}
\label{app:cheatsheet}

\begin{table}[htbp]
  \centering
  \small
  \caption{Class options, given in the square brackets of
  \texttt{\textbackslash documentclass[...]\{boogonreport\}}.}
  \label{tab:options}
  \begin{tabular}{@{}ll@{}}
    \toprule
    Option & Effect \\
    \midrule
    \texttt{zh}          & Chinese: font pipeline, captions and references (XeLaTeX) \\
    \texttt{logo} / \texttt{nologo} & show or hide the brand band \\
    \texttt{times}       & switch body and math to Times \\
    \texttt{numeric}     & numeric citations \texttt{[1]} \\
    \texttt{nonumbering} & unnumbered sections, for white papers \\
    \texttt{leftalign}   & left-aligned title block \\
    \texttt{draft}       & line numbers and visible \verb|\todo{}| \\
    \texttt{fandol}      & portable Chinese character set \\
    \texttt{winfonts}    & Windows system fonts for Chinese \\
    \bottomrule
  \end{tabular}
\end{table}

\begin{table}[htbp]
  \centering
  \small
  \caption{Common commands and environments.}
  \label{tab:commands}
  \begin{tabular}{@{}ll@{}}
    \toprule
    Command & Purpose \\
    \midrule
    \verb|\reportsetup{...}| & title, author, number, version, links, wordmark \\
    \verb|\keywords{...}|    & keywords; works inside or after the abstract \\
    \verb|\secref{...}|      & ``section 3''; idiomatic Chinese phrasing in zh mode \\
    \verb|\cref{...}|        & figure and table references \\
    \verb|\appref{...}|      & appendix references, ``appendix A'' \\
    \verb|\best{}| / \verb|\second{}| & best in bold / second underlined \\
    \texttt{contributions}   & compact list of contributions \\
    \verb|\todo{}| / \verb|\fixme{}| & draft markers, visible only in \texttt{draft} \\
    \bottomrule
  \end{tabular}
\end{table}

\section{Authors and Acknowledgements}
\label{app:authors}

\textbf{Maintainer.} BoogonClothman designs and maintains the template, including the
class, the examples, the documentation and the build pipeline.

\textbf{Acknowledgements.} The layout language owes a debt to the long practice of public
technical reports, in particular the way the DeepSeek and Qwen reports put one result on
the first page. The prose rules follow Bringhurst and the \LaTeX{} community's
conventions.\citep{bringhurst2013elements,mittelbach2023companion}

\textbf{License.} The repository is released under the MIT license: use it, change it,
ship it. Fonts travel with TeX Live and are embedded under their own LPPL/GPL-style
licenses.

\textbf{Source.}\ \url{https://github.com/BoogonClothman/boogon-techreport-kit}


\end{document}
